\chapter{最后要看的第一章作业}

\section{作业1.1 方程的导出与定解条件}
P10三题
\begin{homework}
[1]. 细杆（或弹簧）受某种外界原因而产生纵振动，以 $u(x,t)$ 表示静止时在 $x$ 点处的点在时刻 $t$ 离开原来位置的偏移。假设振动过程中所产生的张力服从胡克定律，试证明 $u(x,t)$ 满足方程
\[
\frac{\partial}{\partial t}\left(\rho(x)\frac{\partial u}{\partial t}\right)
= \frac{\partial}{\partial x}\left(E\frac{\partial u}{\partial x}\right),
\]
其中 $\rho$ 为杆的密度，$E$ 为杨氏模量。
\end{homework}

\begin{homework}
[6] 绝对柔软而均匀的弦线有一端固定，在它本身重力的作用下，此线处于铅垂的平衡位置，试导出此线的微小横振动方程。
\end{homework}

\begin{homework}
[7] 一柔软均匀的细弦，一端固定，另一端是弹性支承。设该弦在阻力与速度成正比的介质中作微小的横振动，试写出弦的位移所满足的定解问题。
\end{homework}



\section{作业1.2 达朗贝尔公式、波的传播}
P19:5题
\begin{homework}[1]
证明方程
\[
\frac{\partial}{\partial x}\!\left[\Bigl(1-\frac{x}{h}\Bigr)^{2}\frac{\partial u}{\partial x}\right]
= \frac{1}{a^{2}}\Bigl(1-\frac{x}{h}\Bigr)^{2}\frac{\partial^{2}u}{\partial t^{2}}
\]
的通解可以写成
\[
u = \frac{F(x-at)+G(x+at)}{h-x},
\]
并求其满足初始条件
\[
u(x,0)=\varphi(x),\qquad u_t(x,0)=\psi(x)
\]
的解。
\end{homework}
\begin{dxtips}
    首先设w=u/g最终我们要把这个式子转换成a²wxx=wtt。gx=-1再复合求导的时候要带入再算。一定要一点一点算。记者如果式子不是用x表示的就先对自变量求导然后再乘自变量对x的导数。
\end{dxtips}
\begin{homework}
[3] 利用传播波法，求解波动方程的古尔萨（Goursat）问题
\[
\begin{cases}
u_{tt}=a^{2}u_{xx}, & -at<x<at,\ t>0,\\[4pt]
u|_{x-at=0}=\varphi(x),\\[4pt]
u|_{x+at=0}=\psi(x),\qquad \varphi(0)=\psi(0).
\end{cases}
\]
\end{homework}
\begin{dxtips}
    计算技巧与目的就是把他写成通解\(F(x-at)+G(x+at)\),把\(F\)和\(G\)换成\(\varphi\)和\(\psi\).带入如何变量替换即可。第二小问还需要再写
\end{dxtips}

\begin{homework}
[5] 求解
\[
\begin{cases}
u_{tt}-a^{2}u_{xx}=0,\ x>0,\ t>0,\\[4pt]
u(x,0)=\varphi(x),\\[4pt]
u_t(x,0)=0,\\[4pt]
(u_x-ku_t)|_{x=0}=0,
\end{cases}
\]
其中 $k>0$。
\end{homework}
\begin{dxtips}
   先正常求解，求出\(F\)和\(G\)和\(\varphi\)的关系记住常数\(C=0\)，第二步也就是为了解决\(x-at\)不能小于\(0\)的问题把边界条件带入运算。注意1，要用到在初值时\(F'()\)和\(G'()\)导数相同然后替换。设系数是\(r\)然后积分由于\(F'(-x)\)里面有\(-\)号所以积分完了就是\(F(-x)=F(x)+A\)
通过\(F\)与\(\varphi\)导关系算出\(A\)然后
\(F(x)=F(-x)+A\)带入
\end{dxtips}
\begin{homework}
[7] 求边值问题
\[
\begin{cases}
u_{tt}-u_{xx}=0,\quad f(t)<x<t,\ t>0,\\[4pt]
u|_{x=t}=\varphi(t),\\[4pt]
u|_{x=f(t)}=\psi(t),
\end{cases}
\]
其中 $\varphi(0)=\psi(0)$，且 $x=f(t)$ 是满足 $|f'(t)|\neq1$ 的光滑曲线。
\end{homework}

\begin{homework}
[8] 求解初值问题
\[
\begin{cases}
u_{tt}-u_{xx}=t\sin x,\qquad t>0,\ -\infty<x<\infty,\\[4pt]
u(x,0)=0,\\[4pt]
u_t(x,0)=\sin x.
\end{cases}
\]
\end{homework}



\section{分离变量法}
P31:4
\begin{dxtips}
    \[
A_n=\frac{2h}{l^2}\cdot \frac{(-1)^{n-1}}{\alpha^{2}}
=\frac{2h}{l^2}\cdot \frac{4l^2(-1)^{n-1}}{(2n-1)^2\pi^2}
=\boxed{\;\frac{8h(-1)^{n-1}}{(2n-1)^2\pi^2}\;},
\qquad
\alpha=\frac{(2n-1)\pi}{2l}.
\]
\end{dxtips}
\begin{dxtips}
    \begin{itemize}
  \item \textbf{$\lambda>0$：}设 $\lambda=k^{2}$，
  \[
  X(x)=A\cos kx+B\sin kx,\qquad T(t)=C e^{-a^{2}k^{2}t}.
  \]

  \item \textbf{$\lambda=0$：}
  \[
  X(x)=A+Bx,\qquad T(t)=C.
  \]

  \item \textbf{$\lambda<0$：}设 $\lambda=-k^{2}$，
  \[
  X(x)=A e^{kx}+B e^{-kx}\ \ (\text{或 }A\cosh kx+B\sinh kx),\qquad
  T(t)=C e^{a^{2}k^{2}t}.
  \]
\end{itemize}
\end{dxtips}
\begin{homework}[1(2)]
求解初边值问题（$0<x<l,\ t>0$）：
\[
\begin{cases}
u_{tt}=a^{2}u_{xx},\\[3pt]
u(0,t)=0,\quad u_x(l,t)=0,\\[3pt]
u(x,0)=\dfrac{h}{l}x,\quad u_t(x,0)=0.
\end{cases}
\]
\end{homework}
\begin{dxtips}
    这题先正常算。最后算系数An
    吧u（x，0）条件带入$$b_n=\frac{2}{l}\int_{0}^{l} f(x)\sin\frac{n\pi x}{l}\,dx.$$就结束了别算了
\end{dxtips}
\begin{homework}[2]
求解初边值问题：
\[
\begin{cases}
u_{tt}=a^{2}u_{xx},\\[3pt]
u(0,t)=0,\quad u(l,t)=A\sin^{2}(\omega t),\\[3pt]
u(x,0)=0,\quad u_t(x,0)=0.
\end{cases}
\]
\end{homework}

\begin{homework}
[4]求解
\[
\begin{cases}
u_{tt}-a^{2}u_{xx}=g,\\[3pt]
u(0,t)=0,\quad u_x(l,t)=0,\\[3pt]
u(x,0)=0,\quad u_t(x,0)=\sin\frac{\pi x}{2l}.
\end{cases}
\]
\end{homework}

\begin{homework}
[6] 求解阻尼波方程
\[
\begin{cases}
u_{tt}+2bu_t=a^{2}u_{xx},\quad b>0,\\[3pt]
u(0,t)=u(l,t)=0,\\[3pt]
u(x,0)=\frac{h}{l}x,\quad u_t(x,0)=0.
\end{cases}
\]
\end{homework}


% 设 \varphi(x,y,z)=x^{3}+y^{2}z，并令 (x',y',z')=(x,y,z)+w
\[
\varphi(x+w_1,y+w_2,z+w_3)
=(x+w_1)^3+(y+w_2)^2(z+w_3).
\]

% 先把它展开
\[
(x+w_1)^3=x^3+3x^2w_1+3xw_1^2+w_1^3,
\]
\[
(y+w_2)^2(z+w_3)=(y^2+2yw_2+w_2^2)(z+w_3)
= y^2z+y^2w_3+2yzw_2+2yw_2w_3+zw_2^2+w_2^2w_3.
\]

% 在球面 |w|=r 上用对称性积分：所有“含 w 的奇次项”积分为 0，
% 以及混合项 w_i w_j (i\ne j) 也为 0；只需用
\[
\iint_{|w|=r} w_i\,dS=0,\qquad
\iint_{|w|=r} w_iw_j\,dS=0\ (i\ne j),\qquad
\iint_{|w|=r} w_i^2\,dS=\frac{4\pi r^4}{3}.
\]

% 因而留下来的只有：常数项 + w_1^2 项 + w_2^2 项
\[
\iint_{|w|=r}\varphi(x+w_1,y+w_2,z+w_3)\,dS
=4\pi r^2(x^3+y^2z)+3x\iint_{|w|=r}w_1^2\,dS
+z\iint_{|w|=r}w_2^2\,dS.
\]
代入上式中的球面积分值，得
\[
\iint_{|w|=r}\varphi\,dS
=4\pi r^2(x^3+y^2z)+4\pi r^4 x+\frac{4\pi}{3}r^4 z.
\]

% 若原式是 \iint_{|w|=r} \frac{\varphi}{r}\,dS，则再除以 r：
\[
\iint_{|w|=r}\frac{\varphi}{r}\,dS
=\frac{1}{r}\iint_{|w|=r}\varphi\,dS
=4\pi r(x^3+y^2z)+4\pi r^3x+\frac{4\pi}{3}r^3z.
\]
\section{4 高维柯西问题}
P64：4
\begin{homework}
[1]利用泊松公式求解三维波动方程的柯西问题：
\[
\begin{cases}
u_{tt}=a^{2}(u_{xx}+u_{yy}+u_{zz}),\\
u(x,y,z,0)=0,\quad u_t(x,y,z,0)=x^{2}+yz.
\end{cases}
\]
以及
\[
\begin{cases}
u_{tt}=a^{2}(u_{xx}+u_{yy}+u_{zz}),\\
u(x,y,z,0)=x^{3}+y^{2}z,\quad u_t(x,y,z,0)=0.
\end{cases}
\]
\end{homework}
\begin{dxtips}
    这里写的是第二题先用标准柏松公式，然后对 $M=x+r\sin\theta\cos\varphi,\ y+r\sin\theta\sin\varphi,\ z+r\cos\theta$。记住 $x,y$ 都有的是 $0$ 到 $2\pi$，$z$ 自己有的是 $0$ 到 $\pi$。$r$ 是一个定值 $at$ 最后带入。$ds=rd\theta* r\sin\theta d\varphi$
为了方便计算写成 $x+w_1,y+w_2,z+w_3$。因为对称性，所以凡事不都是偶次平方的项积分都是 $0$，$w_1$ 平方的二次积分是 $\frac{4}{3}\pi r^2$
一步一步拆开算别忘了最后还要对t求一次导把r代回at
\end{dxtips}
\begin{homework}[2]
试用降维法导出弦振动方程的达朗贝尔公式。
\end{homework}
\begin{dxtips}
先写出标准柏松公式，然后由于是一维的只用设 $z$ 就想 $z=z+r\cos\theta$，$ds$ 什么的不变然后积分因为和 $\varphi$ 无关了直接积分完了把 $2\pi$ 提出来，然后把 $z+r\cos\theta$ 进行变量替换换成 $\xi$ 记得上下变量都要换然后对 $t$ 求导记得里面有一个 $z-at$ 所以还要乘一个 $-1$。记得还有 $\psi$ 的事情呢。
\end{dxtips}
\begin{homework}[5]
求解二维 Klein\text{--}Gordon 型柯西问题：
\[
\begin{cases}
u_{tt}=a^{2}(u_{xx}+u_{yy})+c^{2}u,\\
u(x,y,0)=\varphi(x,y),\\
u_t(x,y,0)=\psi(x,y).
\end{cases}
\]
提示：令 $v=e^{(c/a)z}u$ 并利用三维波动方程。
\end{homework}

\begin{homework}[6]
试用齐次化原理导出非齐次二维波动方程
\[
u_{tt}=a^{2}(u_{xx}+u_{yy})+f(x,y,t)
\]
在初始条件 $u(x,y,0)=0,\ u_t(x,y,0)=0$ 下的求解公式。
\end{homework}
\section{5}
P53
\begin{homework}[1]
试说明：对一维波动方程所描述的波的传播过程一般具有后效现象。
\end{homework}

\begin{homework}[2]
试说明：对一维波动方程，即使初始资料具有紧支集，当 $t\to\infty$ 时其柯西问题的解没有衰减性。
\end{homework}

\begin{homework}[3]
设 $u$ 为初始资料 $\varphi$ 及 $\psi$ 具有紧支集的二维波动方程的解。试证明：对任意固定的 $(x_0,y_0)\in\mathbb{R}^2$，$u(x_0,y_0,t)$ 以 $t^{-1/2}$ 的速度趋于零。
\end{homework}

\section{6}
P62
\begin{homework}[1]
对受摩擦力作用且具有固定端点的有界弦振动，满足方程
\[
u_{tt}=a^{2}u_{xx}-c u_t,
\]
其中常数 $c>0$，证明其能量是减少的，并由此证明方程
\[
u_{tt}=a^{2}u_{xx}-c u_t+f
\]
的初边值问题解的唯一性。
\end{homework}
\begin{dxtips}
    有固定断点两端的速度和位移就都是$0$，先列出$E(t)$公式也就是总能量公式，$dx$前面自然是一次积分从$0$到$l$，$u_t^2$始终不变后面是$a^2u_x^2$，二维就是$u_x^2+u_y^2$后面自然是$ds$前面自然是在$\Omega$区域上积分。然对$t$求导也就是看$E$能量对时间变化率得是小于$0$.对时间偏导了$u_t^2=2u_tu_{tt}$了后面的都会出$u_t$最后提出$u_t$后面的就喝题汗$u_{tt}-a^2u_{xx}=-cu_t$联动上了有了负号再加上平方项积分肯定小于$0$.
\end{dxtips}
\begin{dxtips}
   证明唯一性就是设有两个解然后做差，因为他们都是解肯定都符合一个关系，所以他们相减t=0时刻能量就是0，然后能量变化率还是小于等于0的自然他就只能一直是0所以两个就相等。 
\end{dxtips}
\begin{homework}[5]
考虑波动方程的第三类初边值问题
\[
\begin{cases}
u_{tt}-a^{2}(u_{xx}+u_{yy})=0, & t>0,\ (x,y)\in\Omega,\\[4pt]
u|_{t=0}=\varphi(x,y),\qquad u_t|_{t=0}=\psi(x,y),\\[4pt]
\left.\left(\dfrac{\partial u}{\partial n}+\sigma u\right)\right|_{\Gamma}=0,
\end{cases}
\]
其中 $\sigma>0$ 为常数，$\Gamma$ 为 $\Omega$ 的边界，$n$ 为 $\Gamma$ 上的单位外法线向量。
对上述定解问题的解，定义能量函数
\[
E(t)=\iint_{\Omega}[u_t^2 + a^{2}(u_x^2+u_y^2)]\,dxdy
\ +\ a^{2}\int_{\Gamma}\sigma u^{2}\,ds.
\]
试证明 $E(t)\equiv\text{常数}$，并由此证明上述定解问题的唯一性。
\end{homework}
\begin{dxtips}
    \[
\iint_{\Omega}(u_xu_{xt}+u_yu_{yt})\,dxdy
=\iint_{\Omega}\nabla u\cdot\nabla u_t\,dxdy
=-\iint_{\Omega}(\Delta u)\,u_t\,dxdy+\int_{\Gamma}\frac{\partial u}{\partial n}u_t\,ds.
\]
反正他这个a²那个部分要能和前面凑出第一个条件后面和σ凑出第三个条件
\end{dxtips}
\begin{homework}[7]
设 $u(x,t)$ 是 $[0,1]\times[0,\infty)$ 中初边值问题
\[
\begin{cases}
u_{tt}-u_{xx}=0,\\[4pt]
u|_{x=0}=u|_{x=1}=0,\\[4pt]
u|_{t=0}=0,\qquad u_t|_{t=0}=x^{2}(1-x),
\end{cases}
\]
的解。求
\[
\int_{0}^{1}\left[u_t^{2}(x,t)+u_x^{2}(x,t)\right]\,dx.
\]
\end{homework}
\chapter{最后要看的第二章作业}
\section{2 初边值的分离变量法}[一题]
\begin{dxtips}
    \begin{itemize}
  \item \textbf{1）热方程型（分离后）}
  \[
  T'(t)+a^{2}\lambda\,T(t)=0
  \quad\Rightarrow\quad
  T(t)=C\,e^{-a^{2}\lambda t}.
  \]
  若令 $\lambda=k^{2}$（$k\ge 0$），则
  \[
  \lambda=k^{2}>0:\quad T(t)=C\,e^{-a^{2}k^{2}t},\qquad
  \lambda=0:\quad T(t)=C,
  \]
  若令 $\lambda=-k^{2}$（$k>0$），则
  \[
  \lambda=-k^{2}<0:\quad T(t)=C\,e^{+a^{2}k^{2}t}.
  \]

  \item \textbf{2）波方程型（分离后）}
  \[
  T''(t)+a^{2}\lambda\,T(t)=0.
  \]
  若令 $\lambda=k^{2}$（$k\ge 0$），则
  \[
  \lambda=k^{2}>0:\quad T(t)=C_{1}\cos(akt)+C_{2}\sin(akt),
  \qquad
  \lambda=0:\quad T(t)=C_{1}+C_{2}t.
  \]
  若令 $\lambda=-k^{2}$（$k>0$），则可写成指数形式（不写 $\cosh,\sinh$）：
  \[
  \lambda=-k^{2}<0:\quad T(t)=C_{1}e^{akt}+C_{2}e^{-akt}.
  \]
\end{itemize}
\end{dxtips}
\begin{homework}[1]
用分离变量法求下列定解问题的解：
\[
\begin{cases}
\dfrac{\partial u}{\partial t}
= a^{2}\dfrac{\partial^{2}u}{\partial x^{2}},
& t>0,\ 0<x<\pi,\\[6pt]
u(0,t)=0,\quad
\dfrac{\partial u}{\partial x}(\pi,t)=0,
& t>0,\\[6pt]
u(x,0)=f(x),
& 0<x<\pi.
\end{cases}
\]
\end{homework}
\begin{dxtips}
这个和波动方程的分离变量还不一样。因为这里是T'(t)/T(t)=-λ不是T''(t)/T(t)所以说他的解就是  $$\frac{T’(t)}{T(t)}=-\lambda,\ \text{令 }u=T(t)\Rightarrow du=T’(t),dt\ \Rightarrow\ \int \frac{T’(t)}{T(t)},dt=\int \frac{1}{u},du=-\lambda\int dt\ \Rightarrow\ \ln|T(t)|=-\lambda t+C.$$
$$\ln|T(t)|-\ln|T(0)|=-\lambda t\ \Rightarrow\ T(t)=T(0)e^{-\lambda t}.$$  $\int dt=t+c$
$T(t)=Ce^{-\lambda t}$C把常数项也吸收了。


\end{dxtips}
\begin{homework}[2]
用分离变量法求解热传导方程的初边值问题：
\[
\begin{cases}
\dfrac{\partial u}{\partial t}
= \dfrac{\partial^2 u}{\partial x^2}, & (t>0,\,0<x<1),\\[6pt]
u(x,0)=
\begin{cases}
x, & 0<x\le\dfrac{1}{2},\\[4pt]
1-x, & \dfrac{1}{2}<x<1,
\end{cases}\\[10pt]
u(0,t)=u(1,t)=0, & (t>0).
\end{cases}
\]
\end{homework}
\begin{dxtips}
    $$b_n=\frac{2}{l}\int_{0}^{l} f(x)\sin\frac{n\pi x}{l}\,dx.$$l=1
    背$$b_n=\frac{4}{\pi^2 n^2}\sin\left(\frac{n\pi}{2}\right).$$
\end{dxtips}
\begin{dxtips}
    nπ=α\[
b_n=\frac{4}{\pi^2 n^2}\sin(\alpha),\qquad \alpha=\frac{n\pi}{2}.
\]
\end{dxtips}
\section{3柯西问题——三题}
[很有可能填空]
% 约定：对 $f\in L^1(\mathbb R)$ 定义一维傅里叶变换
% \[
%   \mathcal F[f](\omega)=\hat f(\omega)
%   =\int_{-\infty}^{\infty}f(x)e^{-i\omega x}\,dx .
% \]

%------------------ 题 1 ------------------%
\begin{homework}[1]
求下列函数的傅里叶变换：
\begin{enumerate}[(1)]
  \item $e^{-\eta x^2}\quad(\eta>0)$；
  \item $e^{-a|x|}\quad(a>0)$；
  \item $\displaystyle\frac{x}{(a^2+x^2)^k},\qquad
        \frac1{(a^2+x^2)^k}\quad(a>0,\;k\in\mathbb N)$。
\end{enumerate}
\end{homework}
\begin{dxtips}
$$\int_{-\infty}^{\infty} e^{-x^{2}},dx=\sqrt{\pi}.$$
配方把含x的都弄进平方项然后平方里面平移不变。
$$2\pi i\cdot \operatorname{Res}!\left(\frac{e^{i\omega z}}{(z-ia)(z+ia)},,z=ia\right)
=2\pi i\cdot \frac{e^{i\omega ia}}{2ia}
=\frac{\pi}{a}e^{-a\omega}\quad(\omega>0).$$

$$\operatorname{Res}(f,ia)=\frac{1}{(k-1)!}\lim_{z\to ia}\frac{d^{k-1}}{dz^{k-1}}
\Bigl[(z-ia)^{k}f(z)\Bigr].$$
    指数是相加不是相乘。
    $
\mathcal F[xg](\omega)= -\,i\,\frac{d}{d\omega}\hat g(\omega).
$

$
\frac{d}{d\omega}\Bigl(e^{-i\omega x}\Bigr)=-ix\,e^{-i\omega x}.
$这个就是g=1的特殊情况
然后积分就想留数定理
\end{dxtips}
%------------------ 题 2 ------------------%
\begin{homework}[2]
设 $f(x)$ 在 $(-\infty,\infty)$ 绝对可积，证明其傅里叶变换
\[
\hat f(\omega)=\int_{-\infty}^{\infty}f(x)e^{-i\omega x}\,dx
\]
是连续函数。
\end{homework}
\begin{dxtips}
    这个就是想办法证明w趋紧w0时，$\hat{f}(w)-\hat{f}(w_0)$=0两个积分相减为了将极限号从括号外弄到括号内需要证明$f(x)(e^{iwx}-e^{iw_0x})$有界然后把积分号和极限号交换即可
\end{dxtips}


%------------------ 题 3 ------------------%
\begin{homework}[3]
用傅里叶变换求解三维热传导方程的柯西问题
\[
\begin{cases}
\dfrac{\partial u}{\partial t}
=a^2\left(
\dfrac{\partial^2 u}{\partial x^2}
+\dfrac{\partial^2 u}{\partial y^2}
+\dfrac{\partial^2 u}{\partial z^2}
\right),
& t>0,\ (x,y,z)\in\mathbb R^3,\\[6pt]
u(x,y,z,0)=\varphi(x,y,z),
& (x,y,z)\in\mathbb R^3.
\end{cases}
\]
\end{homework}
\begin{dxtips}
   这题不熟：先对 $x,y,z$ 进行傅立叶变换变成一个变量 $\xi$，$\xi$ 的分量是（$\xi_1,\xi_2,\xi_3$）。目的是把第一个式子变成 $\hat{u}$ 对 $t$ 的导数和 $\hat{u}$ 的关系解出 $\hat{u}$,再带入初值解出完整的 $\hat{u}$,再通过热核公式和$\varphi$的卷积把$\hat{u}$傅立叶逆变换成$u$这里用的三维热核心公式
\end{dxtips}
\section{4极值原理——三题}
p93
\begin{homework}[2]
利用证明热传导方程极值原理的方法，证明：满足方程
\[
\frac{\partial^2 u}{\partial x^2}+\frac{\partial^2 u}{\partial y^2}=0
\]
的函数在有界闭区域上的最大值不超过它在边界上的最大值。
\end{homework}


\begin{dxtips}
    先设他的极大值在区域内然后，因为他是极大值一阶导都等于0，二阶导都小于等于0因为两个二阶导加起来等于0所以他们都是0.一个临域内都是最大值，每个最大值都能往外扩展直到边界。
\end{dxtips}


\begin{homework}[3]
设 $u(x,t)$ 为热传导方程
\[
u_t - a^2 u_{xx} - c u = 0,\qquad c>0,
\]
在矩形区域
\[
R=\{(x,t)\mid 0 < x < l,\ 0<t<T\}
\]
中的解。

已知：
\[
|u(0,t)|,\ |u(l,t)| \le M,\qquad t\in[0,T],
\]
\[
|u(x,0)|\le M,\qquad x\in[0,l].
\]

试证：
\[
|u(x,t)|\le M e^{ct},\qquad (x,t)\in R.
\]
\end{homework}
\begin{dxtips}
    最后要证明的是乘$e^{ct}$设V的时候就设v=u除以$e^{ct}$，最终转化成vt=a²vxx的标准式子好用极值原理，v的边界上小于M整个v小于M，把v换回m即可。
\end{dxtips}


\begin{homework}[5]
设 $u(x,t)$ 是区域 $\{0\le x\le \pi,\ 0\le t<\infty\}$ 中问题的经典解：
\[
\begin{cases}
u_t=u_{xx},\\[4pt]
u_x|_{x=0}=u_x|_{x=\pi}=0,\\[4pt]
u|_{t=0}=\varphi(x),
\end{cases}
\]
其中 $\varphi(0)=\varphi'(\pi)=0$。

证明：
\[
\sup_{0<x<\pi}|u(x,1)|\ \le\ \sup_{0<x<\pi}|\varphi(x)|.
\]

\textit{提示：将函数 $u(x,t)$ 关于直线 $x=\pi$ 作偶延拓到 $x\in(\pi,2\pi)$ 中。}
\end{homework}

\begin{dxtips}
   沿拓是为了不让极大值在边界上。证明这个得引入 $\sup|\varphi(x)|$ 我们先用 $C$ 来表示。弄出 $v^{+}$ 和 $v^{-}$ 一个是 $v-C$ 一个是 $-v-C$ 为的就是证明 $v$ 整个小于 $C$ 任意时刻。然后这两个都符合极大值原理的条件，由于是延拓了边界变成内部了极值只能在底边上。$v^{\pm}$ 任意 $x,t$ 都小于 $0$，$v(x,t)$ 在哪里都小于 $C$ 在 $t$ 等于 $1$ 时自然无论 $x$ 是多少都满足条件。
\end{dxtips}



   

\section{12.1 5解的渐进太六题}
95

\begin{homework}[2]
证明：当 $\varphi(x,y)$ 为 $\mathbb R^2$ 上的有界连续函数，且 $\varphi\in L^1(\mathbb R^2)$ 时，二维热传导方程柯西问题的解，当 $t\to\infty$ 时，以 $t^{-1}$ 衰减率趋于零。
\end{homework}
\begin{dxtips}
    写成热核公式和$\varphi$的卷积t确定时G的最大值就是x，y取1的时候，积分放缩$\varphi$有界所以取他的模长/4πa²当C剩下一个所以u小于等于C/t t趋紧无穷远u趋紧于0.
\end
{dxtips}
下面一题同理



\begin{homework}[3]
证明：当 $\varphi(x,y,z)$ 为 $\mathbb R^3$ 上的有界连续函数，且 $\varphi\in L^1(\mathbb R^3)$ 时，三维热传导方程柯西问题的解，当 $t\to\infty$ 时，以 $t^{-3/2}$ 衰减率趋于零。
\end{homework}

105-106

\begin{homework}[1]
设
\[
u(x_1,x_2,\dots,x_n)=f(r),\qquad
r=\sqrt{x_1^2+x_2^2+\cdots+x_n^2},
\]
是 $n$ 维调和函数（即满足方程
\[
\frac{\partial^2 u}{\partial x_1^2}
+\frac{\partial^2 u}{\partial x_2^2}
+\cdots
+\frac{\partial^2 u}{\partial x_n^2}=0),
\]
试证明：
\[
f(r)=c_1+\frac{c_2}{r^{\,n-2}}\quad (n\ne 2),\qquad
f(r)=c_1+c_2\ln\frac1r\quad (n=2),
\]
其中 $c_1,c_2$ 为任意常数。
\end{homework}

\begin{dxtips}
 1.会求二次偏导
 2.会解常微分方程，常微分方程要用两次变量替换一个是降维把f（x）’=g（x）第二积分的时候变量替换把g（x)'挪到积分号后面。不管是先让r²然后再求导还是什么的，反正如果不会就记住如果求导对象里面没有奇就先对自变量求导再拿自变量对要求导导东西求导这题需要重新计算。
\end{dxtips}
    





\begin{homework}[2]
证明：拉普拉斯算子在球面坐标 $(r,\theta,\varphi)$ 下可以写成
\[
\Delta u
=\frac1{r^2}\frac{\partial}{\partial r}\!\left(r^2\frac{\partial u}{\partial r}\right)
+\frac1{r^2\sin\theta}\frac{\partial}{\partial\theta}\!\left(\sin\theta\frac{\partial u}{\partial\theta}\right)
+\frac1{r^2\sin^2\theta}\frac{\partial^2 u}{\partial\varphi^2}.
\]
\end{homework}
\begin{dxtips}
    $$ds^{2}=h_{1}^{2}dq_{1}^{2}+h_{2}^{2}dq_{2}^{2}+h_{3}^{2}dq_{3}^{2}.$$

$$\Delta u=\frac{1}{h_{1}h_{2}h_{3}}\sum_{i=1}^{3}\frac{\partial}{\partial q_{i}}
\left(\frac{h_{1}h_{2}h_{3}}{h_{i}^{2}}\frac{\partial u}{\partial q_{i}}\right).$$

$$ds^{2}=dr^{2}+r^{2}d\theta^{2}+r^{2}\sin^{2}\theta\,d\varphi^{2},\qquad
h_{r}=1,\ h_{\theta}=r,\ h_{\varphi}=r\sin\theta.$$

$$\frac{h_{r}h_{\theta}h_{\varphi}}{h_{r}^{2}}=r^{2}\sin\theta,\qquad
\frac{h_{r}h_{\theta}h_{\varphi}}{h_{\theta}^{2}}=\sin\theta,\qquad
\frac{h_{r}h_{\theta}h_{\varphi}}{h_{\varphi}^{2}}=\frac{1}{\sin\theta}.$$

$$\Delta u=\frac{1}{r^{2}\sin\theta}\left[
\frac{\partial}{\partial r}\!\left(r^{2}\sin\theta\,u_{r}\right)
+\frac{\partial}{\partial\theta}\!\left(\sin\theta\,u_{\theta}\right)
+\frac{\partial}{\partial\varphi}\!\left(\frac{1}{\sin\theta}\,u_{\varphi}\right)
\right].$$
\end{dxtips}

\begin{homework}[3]
证明：拉普拉斯算子在柱面坐标 $(r,\theta,z)$ 下可以写成
\[
\Delta u
=\frac1r\frac{\partial}{\partial r}\!\left(r\frac{\partial u}{\partial r}\right)
+\frac1{r^2}\frac{\partial^2 u}{\partial\theta^2}
+\frac{\partial^2 u}{\partial z^2}.
\]
\end{homework}
\begin{dxtips}
    $$(h_1,h_2,h_3)=(h_r,h_\theta,h_\varphi)=(1,\ r,\ r\sin\theta).$$

$$(h_1,h_2,h_3)=(h_\rho,h_\varphi,h_z)=(1,\ \rho,\ 1).$$

$$h_r=1\to h_\rho=1,\qquad h_\theta=r\to h_\varphi=\rho,\qquad h_\varphi=r\sin\theta\to h_z=1.$$
\end{dxtips}

\begin{homework}[6]
用分离变量法求解由下述调和方程的第一边值问题所描述的矩形平板
\[
(0\le x\le a,\;0\le y\le b)
\]
上的稳态温度分布：
\[
\begin{cases}
\dfrac{\partial^2 u}{\partial x^2}
+\dfrac{\partial^2 u}{\partial y^2}=0,\\[0.6em]
u(0,y)=u(a,y)=0,\\[0.4em]
u(x,0)=\sin\dfrac{\pi x}{a},\quad u(x,b)=0.
\end{cases}
\]
\end{homework}


\section{6格林公式四题}
\begin{homework}[1]
证明：(2.7) 式对于 $M_0$ 在 $\Omega$ 外与 $\Gamma$ 上的情形成立。
\end{homework}
\begin{equation}\tag{2.7}\label{eq:2.7}
 - \iint_{\Gamma}
 \left[
   u\,\frac{\partial}{\partial n}\!\left(\frac{1}{r}\right)
   - \frac{1}{r}\,\frac{\partial u}{\partial n}
 \right]\mathrm{d}S
 =
 \begin{cases}
   0, & \text{若 } M_0 \text{ 在 } \Omega \text{ 外},\\[4pt]
   2\pi u(M_0), & \text{若 } M_0 \text{ 在 } \Gamma \text{ 上},\\[4pt]
   4\pi u(M_0), & \text{若 } M_0 \text{ 在 } \Omega \text{ 内}.
 \end{cases}
\end{equation}
\begin{dxtips}
    先用一个半径ρ球包住M0，是Bρ，我们要的就是让体积分转面积分的体必须是没有奇点的保证他积分完了以后是0，右边从原来的伽马变成现在的伽马ρ和σρ，σρ就是被挖掉的球都边界，伽马ρ是伽马剩下来的。当ρ趋紧域0伽马ρ就是伽马。考虑到σρ上vdu/dn内项因为u对n的导数有界只有一个ρ抵抗不了2πρ所以这项会一起消失。然后另一项ρ趋紧于0u(M)被固定到一点u（M0）因为对v求n的偏导是-1/r²与2派ρ²的ρ抵消掉所以这一项最后被保留
\end{dxtips}


\begin{homework}[2]
若函数 $u(x,y)$ 是\textcolor{purple}{单位圆}上的调和函数，又已知它在\textcolor{purple}{单位圆周上}的数值为
\[
u = \sin\theta,
\]
其中 $\theta$ 表示极角，问函数 $u$ 在原点之值等于多少？
\end{homework}
\begin{dxtips}[填空]
  \[
u(0,0)=\frac{1}{2\pi}\int_{0}^{2\pi}u(1,\theta)\,d\theta
=\frac{1}{2\pi}\int_{0}^{2\pi}\sin\theta\,d\theta=0.
\]
\end{dxtips}

\begin{homework}[3]
设 $u(M)$ 在 $\Omega$ 内调和，$M_0$ 是 $\Omega$ 中的任意一点。
$B_a$ 是以 $M_0$ 为球心、$a$ 为半径的球体，其体积为 $|B_a|$，证明：
\[
u(M_0)=\frac{1}{|B_a|}\iiint_{B_a} u\,dV.
\]
\end{homework}

\begin{dxtips}[老师提到过]
   先一通代写出2.7情况3然后把r=a带入还是有du/dn内一项也就是没有1/a²内一项，再用一次格林公式当v=1又转会体积分由于调和积分就是0.然后对这个结果把a换会r然后再从0到a积分一次就得到结果 
\end{dxtips}


\begin{homework}[4]
证明：当 $u(M)$ 在闭曲面 $\Gamma$ 的外部调和，并且在无穷远处成立
\[
u(M)=O\!\left(\frac{1}{r_{OM}}\right),\qquad
\frac{\partial u}{\partial r}=O\!\left(\frac{1}{r_{OM}^2}\right)
\quad (r_{OM}\to\infty),
\]
而 $M_0$ 是 $\Gamma$ 外的任一点，则公式 (2.6) 仍成立。
\begin{equation}\tag{2.6}\label{eq:2.6}
u(M_0)
= -\frac{1}{4\pi}
  \iint_{\Gamma}
  \left[
    u(M)\,\frac{\partial}{\partial n}
      \!\left(\frac{1}{r_{M_0M}}\right)
    - \frac{1}{r_{M_0M}}\,
      \frac{\partial u(M)}{\partial n}
  \right]\mathrm dS_M .
\end{equation}
\end{homework}
\begin{dxtips}
    先画一个大球包含伽马和M0然后大球挖掉M0领域和伽马，在这个上面他调和，然后边界被三块，伽马，SR大球表面和SρM0表面。M0部分照旧算出4πu(M0),大球面上SR用题目给的条件对积分内容一起放缩结果就是他们加在一块都是小于C/R三次方级别就算乘表面积还是底下多一个R，这个R是要趋紧域无穷的。
\end{dxtips}


\section{7}
P121
\begin{homework}[1]
证明格林函数的性质 3 及性质 5。
\end{homework}
\begin{dxtips}[性质3]
    切掉奇点，球就有内表面和外表面外表面=0整个大于等于0，内表面半径越小1/4πr越大，g有最大值A 总能让1/4πr-A大于0内表面也大于0所以内部大于0.然后证明G小于1/4πr也就是g。用F=1/4πr-G得到g，g在Γ上等于1/4πr大于0，g在omiga内布也就大于0所以F大于0所以G小于1/4πr
\end{dxtips}
\begin{dxtips}[性质5]
求的是dG/dn的积分在伽马上积分等于-1，还是挖掉M0然后把边界分成两个部分M0的表面和剩下的，ε趋紧于0g内项趋紧于0就剩下对1/4πr求导再积分，外法向和内法向相反所以最终积分乘表面积4πr²全部消掉
    
\end{dxtips}

\begin{homework}[5]
求半圆区域上狄利克雷问题的格林函数。
\end{homework}
\begin{dxtips}
  $$
G(M,M_0)=\frac{1}{2\pi}\left(\ln\frac{1}{r_{M_0M}}-\ln\frac{R}{\rho_0,r_{M_1M}}\right).
$$
就是 $r_{M_0M}-r_{M_1M}$ 就是抵消圆周上的。比例上 $R/\rho_0$。$M_0$ 就是 $(\rho_0,\varphi)$，$M_1$ 是 $(\rho_1,\varphi)$ 角度相同长度不同。$M_2$ 是 $M_0$ 的反点都去负号，$M_3$ 是 $M_1$ 点反衬点都取负号。然后因为算的半径都是与 $M$ 的距离所以 $M$ 是 $(\rho,\theta)$。用极坐标求长度公式。
\end{dxtips}
\begin{dxtips}
    若用极坐标 \(A=(\rho_a,\alpha),\,B=(\rho_b,\beta)\)，则等价写成（余弦定理）：
\[
|A-B|
=\sqrt{\rho_a^{2}+\rho_b^{2}-2\rho_a\rho_b\cos(\alpha-\beta)}.
\]
\end{dxtips}